
\chapter{Technical Analysis of MATLAB Script: \texttt{linkplatonic.m}}
\date{}



\maketitle

\section{Overview}

The \texttt{linkplatonic.m} file is a development script designed to construct and visualize 3D polyhedral skeletons using complex plane mathematics and spherical projection techniques.

\begin{itemize}
    \item \textbf{Main Purpose}: To define vertex coordinates and edge connectivity for polyhedra and map them from a 2D plane to a 3D sphere.
    \item \textbf{Type of Analysis}: The script performs coordinate transformations, numerical root-finding for geometric optimization, and topological mapping.
    \item \textbf{Mathematical Context}: The code is grounded in \textbf{Graph Theory} (nodes and adjacency permutations) and \textbf{Stereographic Projection} (the inverse mapping from $\mathbb{C} \to S^2$).
\end{itemize}

\section{Step-by-Step Code Analysis}

\subsection{Variable Initialization and Trial Selection}
\textbf{Explanation}: The script uses a \texttt{switch} statement controlled by the variable \texttt{Trial} to select between different geometric construction methods.

\begin{lstlisting}
Trial=3;  % Selection of the active geometric trial
switch Trial
\end{lstlisting}

\subsection{Section 1: Tetrahedron Construction (Trial 1)}
\textbf{Explanation}: This section calculates the geometry of a tetrahedron as a pyramid with $N$ base edges. It uses \texttt{fzero} to solve a transcendental equation to find the vertical offset $z$ that ensures geometric symmetry upon projection.

\begin{lstlisting}
case 1 
N=6; 
theta=@(z) acos(z) - 2 / N * pi * sqrt (1 - z .^ 2);
z=-fzero(theta,[-0.7,.9]);
r=sqrt(1-z^2);
pts=[r*exp((1:N)/N*2*pi*1i) , 0].';
phi_base=[N 1:N-1];
idA=(1:4*N)';
phi=zeros(N+3*N,1);
s=zeros(N+3*N,1);
phi(1:N)=[2:N 1];
s(1:N)=(1:N)';
for i=1:N
    phi(i+[N 2*N 3*N])=i+[ 2*N 3*N N];
    s(i+[N 2*N 3*N])=[ i N+1  phi_base(i)];
end
pos=pts(s);
V=plane2sphere(pos);
V(:,3)=-V(:,3);
linkdraw(V,phi,idA,.25,0)
\end{lstlisting}

\subsection{Section 2: N-gonal Pyramid Mapping (Trial 2)}
\textbf{Explanation}: This block constructs a pyramid structure. It centers the vertices in the complex plane, normalizes them, and generates both a 2D planar graph and a 3D spherical projection.

\begin{lstlisting}
case 2 
N=6;
idA=(1:N)';
phi=[2:N 1]'; pos=exp(idA/N*2*pi*1i);
idA_=[idA; idA+N; idA+2*N; idA+3*N ];
phi_=[phi; idA+2*N; idA+3*N; idA+N ];
s_=[idA; phi; idA; repmat(N+1,N,1) ];
pos_=[pos-1/N*sum(pos); 0];
pos_=pos_/mean(abs(pos_(1:end-1)));

subplot(1,2,1)
linkdraw(pos_(s_),phi_,idA_,.25,0);
subplot(1,2,2)
linkdraw(complex2sphere(pos_(s_)),phi_,idA_,.25,0);
\end{lstlisting}

\subsection{Section 3: Cube and Prism Architecture (Trial 3)}
\textbf{Explanation}: The active trial constructs a prism by defining two concentric rings of vertices in the complex plane (one scaled by $0.5$ and another by $2.0$). It uses an inverse adjacency map (\texttt{phinv}) to link the layers, effectively creating the side faces of the solid.

\begin{lstlisting}
case 3 
N=5;
phi=[2:N, 1]';
pos=exp((1:N)/N*2*pi*1i).';
idA=(1:N)';
center=pos-sum(pos)/N;
pos=center/mean(abs(center));

pos_=[1/2*pos; 2*pos];
phinv=zeros(numel(phi),1); phinv(phi)=idA;
idA_=[idA; idA+N; idA+2*N; idA+3*N; idA+4*N; idA+5*N];
phi_=[phi; phinv+N; idA+3*N; idA+4*N; idA+5*N; idA+2*N];
s_=[idA; phi+N; phi+N; phi; idA; idA+N];

subplot(1,2,1)
h1=linkdraw(pos_(s_),phi_,s_,0,0);
set(h1,'showarrowhead','off')
subplot(1,2,2)
h2=linkdraw(complex2sphere(pos_(s_)),phi_,s_,0,0);
set(h2,'showarrowhead','off')
\end{lstlisting}

\section{Expected Outputs and Visuals}

The execution of Trial 3 (default) results in a dual-subplot figure:



\begin{enumerate}
    \item \textbf{Subplot 1 (2D View)}: A Schlegel-like diagram of a pentagonal prism ($N=5$). It shows two concentric pentagons where corresponding vertices are connected by links.
    \item \textbf{Subplot 2 (3D View)}: A wireframe on a sphere. The inner vertices map to the southern hemisphere and the outer vertices map to the northern hemisphere, visualizing the prism "wrapped" around the sphere.
\end{enumerate}
